\section{Conclusion}

In this paper, we have demonstrated that the intrinsic similarity among character components can be leveraged to substantially enhance recognition performance for characters situated in the tail of long-tail distributions. This improvement is especially critical for new or infrequently encountered characters, as their representation is often underemphasized in traditional deep learning models. By recognizing and exploiting these shared structural elements, the network gains an enriched understanding of how individual strokes or radical components form complete characters, ultimately boosting accuracy where it is most needed.

Moreover, we introduced the spindle-shaped network architecture as a powerful means of extracting and refining local character component features. This design strategically expands channel capacity in mid-level layers, facilitating a richer embedding space and promoting more reliable feature representations that capture subtle distinctions between visually similar glyphs. The resultant focus on local components not only bolsters recognition of the rarest characters but also provides a more robust framework for accommodating newly encountered symbols. Through additional analyses, we verified the effectiveness of SpindleNet in adapting to both head and tail classes, ensuring balanced performance across the entire character set.

Extensive experiments conducted on three challenging Chinese ancient book datasets, namely TKH, MTH1000, and MTH1200, further validate the advantages of our method. The outcomes confirm that SpindleNet achieves state-of-the-art performance, surpassing conventional solutions by handling extreme data imbalance and preserving fidelity for high-frequency characters. In challenging historical textual corpora that feature a broad and often sparsely sampled set of characters, such improved recognition capabilities are indispensable for accurate transcription and subsequent scholarly research.

Looking ahead, we plan to extend our framework to a weakly-supervised training paradigm, a direction that has already shown promise in scene text recognition tasks. This approach could reduce the burden of precise annotations while potentially increasing the volume of available training samples. By integrating weak labels or partial annotations, future models may become even more adept at unraveling the intricate forms of rare and unseen characters, thereby advancing the state of the art for ancient document understanding.
